\chapter{Architecture}

\section{Overview}

heapx is organized around a small public dispatch layer and several private backends.
The public layer lives primarily in \texttt{src/heap.c}.
It owns backend selection, common handle validation, null-handling, invariant dispatch, node-pool utilities, and allocation-size helpers.

Each backend lives in \texttt{src/heaps/}.
The backend is responsible for its own structural representation and algorithmic repairs.
This split keeps the public API contract consistent while allowing each heap implementation to use a different internal model.

\section{Common Base Object}

Every concrete backend embeds \texttt{struct heapx\_heap} as its first field.
This is a common C pattern for implementing a small form of object-oriented dispatch.
Because the base is first, a pointer to the concrete object can be safely exposed as a pointer to the base object, and backend code can cast it back to the concrete type.

\begin{lstlisting}[style=cstyle, caption={Simplified common base object.}]
struct heapx_heap {
    const struct heapx_vtable *vtable;
    heapx_cmp_fn cmp;
    uint64_t id;
    struct heapx_handle_slot *handle_slots;
    size_t handle_slot_count;
    size_t handle_slot_capacity;
    size_t free_handle_slot;
};
\end{lstlisting}

The common base contains only state that is shared by all implementations.
Concrete storage such as arrays, root pointers, node pools, consolidation tables, and rank tables belongs to the backend structs.

\section{Vtable Dispatch}

The internal vtable mirrors the public heap API.
Public functions validate their common preconditions, then call the backend implementation through the vtable.

\begin{lstlisting}[style=cstyle, caption={Simplified vtable shape.}]
struct heapx_vtable {
    void (*destroy)(struct heapx_heap *heap);
    int (*insert)(struct heapx_heap *heap, void *item);
    int (*insert_handle)(
        struct heapx_heap *heap,
        void *item,
        struct heapx_handle *out
    );
    int (*decrease_key)(struct heapx_heap *heap, void *owner);
    void *(*remove)(struct heapx_heap *heap, void *owner);
    int (*contains)(const struct heapx_heap *heap, const void *item);
    void *(*peek_min)(const struct heapx_heap *heap);
    void *(*extract_min)(struct heapx_heap *heap);
    size_t (*size)(const struct heapx_heap *heap);
    int (*empty)(const struct heapx_heap *heap);
    int (*check_invariants)(const struct heapx_heap *heap);
};
\end{lstlisting}

The key design choice is that targeted vtable operations receive an already resolved backend owner pointer.
For example, public \texttt{heapx\_decrease\_key()} resolves the handle once and then calls the backend with the node or locator associated with that handle.
This avoids duplicating public handle validation in every backend.

\section{Factory Function}

\texttt{heapx\_create()} validates the comparator and dispatches construction to the selected backend factory.
Unsupported implementation selectors return \texttt{NULL}.

\begin{lstlisting}[style=cstyle, caption={Factory logic, simplified.}]
switch (implementation) {
case HEAPX_BINARY_HEAP:
    heap = binary_heap_create(cmp);
    break;
case HEAPX_FIBONACCI_HEAP:
    heap = fibonacci_heap_create(cmp);
    break;
case HEAPX_KAPLAN_HEAP:
    heap = kaplan_heap_create(cmp);
    break;
default:
    return NULL;
}
\end{lstlisting}

The factory also enables invariant checks after successful construction when internal checks are enabled.

\section{Public Dispatch Responsibilities}

The public layer owns behavior that should be independent of backend:

\begin{itemize}
    \item rejecting a missing comparator;
    \item accepting \texttt{heapx\_destroy(NULL)};
    \item returning defined values for null heap queries;
    \item resolving handles before targeted operations;
    \item running invariant checks after successful mutations;
    \item releasing the common handle table during destruction.
\end{itemize}

This keeps backend code focused on heap algorithms.
It also makes the API easier to document because edge-case behavior is centralized.

\section{Backend Responsibilities}

Backends own:

\begin{itemize}
    \item concrete storage layout;
    \item insertion and repair algorithms;
    \item remove and extract-min algorithms;
    \item pointer-identity search for \texttt{contains};
    \item backend-specific invariant checking;
    \item backend-owned metadata allocation and release.
\end{itemize}

The separation is especially useful in an experimental repository.
Adding a new backend requires a concrete object, a static vtable, invariant checks, and a factory, while the public API remains unchanged.
